% ==============================================================================
% T0/FFGFT-Theorie: Dokument 242
% Titel: Die Skalenleiter – Bausteine und Felder durch alle Layer der Materie.
%        Eine pädagogische Skizze zu interskalaren Kopplungen in FFGFT.
% Autor: Johann Pascher
% Datum: Mai 2026
% Status: Hypothese, kein Beweis (Plausibilitätsskizze).
% Bezug: Dokumente 159 (Konzertsaal/Modenstruktur), 133 (K_frak und RG-Lauf),
%        167 (LiNbO3), 168 (Periodensystem), 182 (Universum Maximalskala),
%        205 (Narrativ), 230 (Hilbertraum-Brücke), 241 (Zwei Schichten)
% ==============================================================================

\section*{Vorbemerkung}

Dieses Dokument ist eine \emph{Plausibilitätsskizze}, keine Ableitung.
Es enthält Hypothesen, keine Beweise. Wo eine quantitative Aussage
gemacht wird, ist sie als Größenordnungs-Schätzung zu lesen, nicht als
Prädiktion. Die zentrale Frage --- \emph{Wie stark koppeln die
verschiedenen Skalen der Materie geometrisch miteinander?} --- ist eine
offene Forschungsfrage, kein gelöstes Problem.

Der Sinn dieses Dokuments ist nicht, eine Antwort zu liefern, sondern
die Frage \emph{sauber zu stellen} --- so, dass jeder, der die FFGFT
durcharbeitet, die Skalenleiter zwischen Elementarteilchen und
Kosmologie nicht als Lücke übersieht, sondern als Forschungsterrain
erkennt.

Die FFGFT ist bisher \emph{am untersten Layer} ausgearbeitet:
Elementarteilchen, Moden auf dem 4D-Torus, Wicklungszahlen. Die obere
Grenze --- die kosmologische Maximalskala --- ist in Dok.~182 separat
behandelt. Was zwischen beiden liegt --- Atome, Moleküle, Kristalle,
Gase, Flüssigkeiten, Planeten, Sterne, Galaxien --- ist in der
Dokumentenreihe an einzelnen Stellen punktuell adressiert (Dok.~167
für Kristalle, Dok.~168 für Atome), aber nicht als zusammenhängende
Skalenleiter narrativ entwickelt. Genau das versucht dieses Dokument
nachzuholen, im pädagogischen Stil von Dok.~241.

\section{Ein allgemein ungelöstes Problem}

Bevor FFGFT ins Spiel kommt, lohnt ein Blick darauf, was die
Standardphysik zu dieser Frage tatsächlich sagt --- oder nicht sagt.

Das Standardmodell (SM) der Teilchenphysik ist auf seiner Skala
außerordentlich präzise. Es beschreibt Quarks, Leptonen und ihre
Wechselwirkungen mit einer Genauigkeit, die keine andere physikalische
Theorie erreicht. Aber es hat \emph{keine Hierarchietheorie}: Es gibt
keine Ableitung im SM, warum Quarks in Hadronen gebunden sind, Hadronen
in Kernen, Kerne in Atomen --- und warum diese Strukturen gerade so
aussehen, wie sie aussehen. Das ist eine Beobachtung, keine Erklärung.

\subsection*{Was das Standardmodell hat}

Die Standardphysik begegnet diesem Problem durch \emph{effektive
Theorien}: QCD für Quarks und Gluonen; Kernphysik als
phänomenologische Erweiterung; Atomphysik über QED. Die Übergänge
zwischen diesen Ebenen sind nicht abgeleitet --- sie sind empirisch
justiert. Niemand hat aus der QCD-Lagrangedichte allein berechnet,
warum ein Proton die Masse hat, die es hat, geschweige denn, warum ein
Wasserstoffatom gerade \emph{so} aufgebaut ist. Gitterfeldtheorie kann
das Proton inzwischen numerisch simulieren; \emph{ableiten} im Sinne
von verstehen --- warum genau diese Struktur und keine andere --- kann
sie nicht.

Der nächste Verwandte zu einer skalenübergreifenden Erklärung ist die
\emph{Renormierungsgruppe} (RG). Die RG beschreibt, wie
Kopplungskonstanten mit der Energie skalieren. Asymptotische Freiheit
in der QCD ist genau das: Bei hoher Energie werden Quarks quasi-frei,
bei niedriger Energie stark gebunden. Das ist ein echter
Skalierungseffekt. Aber die RG läuft \emph{vertikal} durch
Energieskalen und sagt nichts darüber, wie die Struktur auf einer
Skala die nächsthöhere \emph{formt}. Sie beschreibt den Übergang der
Kopplungen, nicht die Entstehung von Strukturen.

\subsection*{Das eigentliche Leck}

Die eigentliche Frage ist: Warum trägt ein Atom, als Konsequenz der
Quarkstruktur des Protons, gerade die Struktur, die es trägt? Das SM
kann das nicht beantworten. Die Frage, wie die Eigenschaften des
Protons (Masse, Spin, magnetisches Moment) aus der QCD entstehen, ist
bis heute numerisch schwer und konzeptuell offen. Und eine Ebene
höher: Warum bilden Protonen und Neutronen gerade die Kerne, die sie
bilden, mit gerade diesen Bindungsenergien? Die Kernphysik ist
empirisch --- keine Ableitung aus der QCD.

Das ist kein Mangel an Akribie. Es ist ein strukturelles Merkmal der
Standardphysik: Sie beschreibt die Mechanismen \emph{innerhalb} jeder
Ebene sehr gut, und die \emph{lokalen Übergänge} von einer Ebene zur
nächsten durch passend gewählte Kopplungsparameter. Was sie nicht
beschreibt, ist, ob und wie eine Geometrie auf einer Skala
\emph{durch} die Skalenleiter hindurch wirkt --- ob also das, was auf
der Quark-Ebene fundamental ist, in irgendeiner Form in der Struktur
des Atoms oder des Kristalls noch sichtbar ist.

\subsection*{Ein empirischer Beleg: externe Beeinflussung als Strukturbeweis}

Es gibt ein Argument, das keine Theorie benötigt und dennoch direkt auf
den Kern der Frage zeigt.

Wir beeinflussen molekulare Strukturen durch externe Felder: Ein
elektromagnetisches Feld ändert die Konformation eines Proteins; ein
mechanischer Druck verschiebt einen Kristallisationspfad; eine
Temperaturänderung treibt eine chemische Reaktion in eine andere
Richtung. Diese Beeinflussung ist kein Spezialeffekt --- sie ist
experimenteller Alltag in Chemie, Biophysik und Materialwissenschaft.

Aber sie hat eine strukturelle Konsequenz, die leicht übersehen wird:
\emph{Damit ein externes Feld auf einer höheren Skala eine molekulare
Struktur auf einer tieferen Skala verändern kann, müssen
Kopplungskanäle zwischen diesen Skalen physikalisch existieren.}
Diese Kanäle sind keine Eigenschaft des externen Feldes. Sie sind eine
Eigenschaft der Struktur selbst.

Und hier kommt der entscheidende Schritt: Kopplungskanäle, die
physikalisch existieren, sind bidirektional. Wenn eine makroskopische
Feldkonfiguration auf die Elektronenhülle eines Moleküls wirken kann,
dann wirkt umgekehrt die Substruktur des Moleküls --- die Anordnung
der Atome, die Elektronendichte, die Kerngeometrie --- auf die
übergeordnete Struktur. Nicht durch einen externen Eingriff, sondern
als systeminterne Propagation.

Das ist kein spekulativer Schluss. Er zeigt sich bereits in der
Proteinfaltung: Ein Protein faltet sich ohne externen Eingriff in seine
Gleichgewichtsstruktur. Die Information darüber, \emph{wie} es sich
faltet, sitzt in der Aminosäuresequenz --- also auf der chemischen
Skala. Diese Information propagiert aufwärts zur Makromolekül-Geometrie.
Das ist interskalare Kopplung, systemintern, ohne externe Quelle.

Die Standardphysik beschreibt diese Kanäle lokal und korrekt:
Coulomb-Wechselwirkung, van-der-Waals, Wasserstoffbrücken. Was sie
nicht beschreibt, ist, ob diese lokalen Kanäle Ausdruck einer
\emph{durchgehenden geometrischen Struktur} sind --- oder nur eine
Sammlung zufällig passender Mechanismen. FFGFT behauptet Ersteres.
Das SM schweigt dazu.

\subsection*{Stabilität als versteckte Voraussetzung}

Hinter dem Schweigen des SM zur interskalaren Kopplung liegt eine
methodische Entscheidung, die selten explizit gemacht wird: Das SM
behandelt stabile übergeordnete Strukturen als \emph{Randbedingungen},
nicht als \emph{Erklärungsziele}. Ein Atom existiert --- gut, dann
rechnen wir mit ihm. Ein Kristall existiert --- gut, dann beschreiben
wir seine Bandstruktur. Warum diese Strukturen überhaupt stabil sind,
und warum sie gerade diese Form haben und keine andere, wird nicht
gefragt. Stabilität wird als gegeben hingenommen.

Das versteckt eine fundamentale Annahme: dass die interskalaren Kräfte,
die diese Stabilität \emph{erzeugen und aufrechterhalten}, irgendwie
selbstverständlich sind. Ein Atom ist stabil, weil die
elektromagnetische Kraft zwischen Kern und Elektronen eine bestimmte
Stärke hat und weil die Quantenmechanik bestimmte Orbitale
stabilisiert. Das ist eine präzise Balance zwischen Kräften
verschiedener Skalen --- Kernkräfte auf der Femtometerskala,
elektromagnetische auf der Ångströmskala. Diese Balance ist nicht
trivial. Sie ist das Ergebnis genau jener interskalaren Kopplung, die
das SM nicht erklärt. Das SM \emph{nutzt} die Stabilität, die aus
dieser Kopplung entsteht, ohne sie zu begründen.

\subsection*{Das Planetenbahn-Argument}

In der Himmelsmechanik ist es selbstverständlich, dass eine
Planetenbahn \emph{nie} ein reines Zweikörperproblem ist. Jupiter
stört Saturn. Der Mond stört die Erdbahn. Die Gezeiten koppeln
Rotation und Bahnentwicklung. Kleine Kopplungen akkumulieren sich über
Jahrmillionen zu strukturbildenden Effekten --- Bahnresonanzen, Laplace-
Resonanzen bei Jupiters Monden, die Stabilisierung der Erdachsenneigung
durch den Mond. Niemand in der Himmelsmechanik sagt: ``Wir
vernachlässigen Jupiter, weil er weit weg ist.'' Die Astronomie hat
gelernt, dass auch kleine Kopplungen über lange Zeiten
\emph{strukturkonstitutiv} sind.

Auf der molekularen und atomaren Ebene wird dieselbe Logik nicht
angewendet. Dort gilt die Skalentrennung als Arbeitsprinzip --- QCD
für Quarks, QED für Atome, Festkörperphysik für Kristalle, und zwischen
diesen Theorien liegen sorgfältig gezogene Grenzen. Die Frage, ob
kleine Restkopplungen zwischen diesen Skalen über ausreichend lange
Zeitskalen oder in hinreichend präzisen Messungen strukturell sichtbar
sind, wird gar nicht erst gestellt. Nicht weil die Antwort ``Nein''
lautet, sondern weil die effektiven Theorien so gut funktionieren, dass
niemand nachgefragt hat, was sie stillschweigend wegintegrieren.

Das ist keine physikalische Notwendigkeit, sondern eine historisch
gewachsene Konvention. Die Himmelsmechanik und die Teilchenphysik
behandeln interskalare Kopplungen asymmetrisch --- nicht aus einem
Prinzip heraus, sondern aus unterschiedlichen Traditionen. FFGFT
beansprucht, beide unter derselben geometrischen Logik zu vereinen:
Kopplungen existieren auf allen Skalen, ihre Stärke skaliert
geometrisch mit dem logarithmischen Abstand, und Stabilität übergeordneter
Strukturen ist nicht selbstverständlich, sondern \emph{erklärungs\-bedürftig}.

\subsection*{Warum Astronomen Kopplungen sehen, Teilchenphysiker nicht}

Es gibt einen einfachen Grund, warum interskalare Kopplungen in der
Himmelsmechanik so augenfällig sind und in der Teilchenphysik nicht ---
und er hat nichts damit zu tun, ob die Kopplungen grundsätzlich
vorhanden sind. Er hat damit zu tun, \emph{wie weit auseinander} die
beteiligten Skalen liegen.

In der Himmelsmechanik spielen Jupiter und Saturn auf derselben Bühne.
Sie sind beide Planeten, beide im Sonnensystem, beide auf vergleichbaren
Bahnen. Der Abstand zwischen ihren Skalen ist winzig --- fast null,
wenn man die volle Spannweite der Naturskalen betrachtet. Deshalb ist
die gegenseitige Beeinflussung groß genug, um über Jahrmillionen
sichtbare Spuren zu hinterlassen.

In der Teilchenphysik dagegen liegen die Skalen, zwischen denen man
Kopplungen suchen würde, unvorstellbar weit auseinander. Zwischen einem
Quark und einem Atom passen acht Zehnerpotenzen. Zwischen der
Planck-Skala und der Skala, auf der das Higgs-Teilchen sitzt,
passen siebzehn. Eine Kopplung, die über so viele Größenordnungen
wirkt, wird mit jeder weiteren Größenordnung schwächer --- so schwach,
dass sie mit keinem heutigen Messgerät nachweisbar wäre.

Das ist der eigentliche Grund, warum die Standardphysik mit getrennten
Skalen so gut auskommt: nicht weil die Kopplungen fehlen, sondern weil
die Abstände so gewaltig sind, dass die Kopplungen praktisch
verschwinden. Astronomie und Teilchenphysik widersprechen sich nicht
--- sie schauen auf verschiedene Abschnitte derselben Leiter.

\subsection*{Einfachheit in der Tiefe --- warum Grundbausteine so stabil sind}

Es gibt noch einen zweiten Grund, warum das Bild nach unten immer
stabiler wird. Er betrifft nicht die Abstände zwischen den Skalen,
sondern die Natur dessen, was auf jeder Skala sitzt.

Je tiefer man auf der Skalenleiter vordringt, desto \emph{einfacher}
werden die Strukturen. Ein Elektron hat keine inneren Teile, keine
Unterstruktur, keine Freiheitsgrade, in die es sich umordnen könnte.
Ein Quark ist ähnlich einfach. Diese Einfachheit ist kein Zufall ---
sie ist der Grund für ihre außerordentliche Stabilität. Eine einfache
Struktur kann nicht zerfallen, weil es nichts Einfacheres gibt, in das
sie zerfallen könnte. Der tiefste Zustand auf einer Skala ist stabil
schlicht deshalb, weil darunter nichts mehr ist.

Je höher man auf der Leiter steigt, desto komplexer werden die
Strukturen. Ein Atom hat schon viele Freiheitsgrade --- Elektronen auf
verschiedenen Bahnen, Kernspin, Schwingungsmoden. Ein Molekül hat noch
mehr. Ein Protein hat so viele, dass es in Tausenden verschiedener
Konformationen vorliegen kann. Auf der planetaren Skala ist die
Komplexität so gewaltig, dass von einer geometrischen Grundstruktur
praktisch nichts mehr direkt sichtbar ist --- sie ertrinkt im Rauschen
der Milliarden von Atomen und ihrer Wechselwirkungen.

Diese Zunahme der Komplexität nach oben erklärt, warum Stabilität
relativ wird: Ein Atom ist stabil unter normalen Bedingungen, aber
ionisierbar mit genug Energie. Ein Molekül ist stabil bei
Raumtemperatur, aber zersetzbar durch Wärme oder Licht. Die Stabilität
nimmt ab, je weiter man von den einfachen Grundbausteinen entfernt ist.

\subsection*{Die Umkehrung: wie viel Aufwand es kostet, einfache Strukturen zu brechen}

Denselben Sachverhalt kann man von der anderen Seite betrachten ---
und er wird dabei besonders greifbar.

Um ein Elektron aus einem Wasserstoffatom herauszulösen, genügt
ultraviolettes Licht. Das ist Energie, die wir mit einer gewöhnlichen
Lampe erzeugen können. Um einen Atomkern zu spalten, braucht man
Neutronenbeschuss oder extreme Temperaturen --- Millionen Grad, wie sie
in Sternen herrschen oder in einer Wasserstoffbombe erzeugt werden.
Das ist eine Million Mal mehr Energie. Um schließlich ein Proton in
seine Quarks zu zerlegen, braucht man einen der größten
Teilchenbeschleuniger der Welt --- und selbst dann sieht man die
Quarks nicht frei, weil sie sofort wieder neue Verbindungen eingehen.

Diese Staffelung des Aufwands ist kein Zufall. Sie zeigt direkt, wie
tief eine Struktur auf der Skalenleiter sitzt. Je tiefer, desto
einfacher --- und desto mehr Energie braucht man, um sie aufzubrechen,
weil es nichts gibt, in das sie ``fallen'' könnte. Die Energieschwelle
ist gleichsam der experimentelle Fingerabdruck der Tiefe.

Umgekehrt gelesen: Dass die Kernspaltung Millionen Mal mehr Aufwand
erfordert als die Ionisierung eines Atoms, ist der direkte Beweis
dafür, dass Atomkerne auf einer tieferen, einfacheren Ebene sitzen als
Elektronen. Und dass selbst das noch nicht die unterste Ebene ist,
zeigt der Aufwand, der nötig wäre, um Quarks zu isolieren. Die
Skalenleiter ist keine Metapher --- sie ist in diesen Energieschwellen
physikalisch eingeschrieben.

Das Standardmodell beschreibt diese Schwellen präzise und korrekt.
Was es nicht erklärt, ist, warum die Abstände zwischen den Skalen
gerade so sind wie sie sind --- warum nicht halb so groß oder doppelt
so groß. FFGFT behauptet, dass diese Abstände keine Zufälle sind,
sondern einer geometrischen Ordnung folgen. Ob das stimmt, ist offen.
Dass die Frage berechtigt ist, zeigt die Staffelung selbst.

\subsection*{Das Hierarchieproblem als Symptom}

Das bekannteste Symptom dieses Lecks ist das \emph{Hierarchieproblem}:
Warum ist die Gravitationsskala (Planck-Masse) so viele
Größenordnungen größer als die elektroschwache Skala (Higgs-Masse)?
Das SM hat keine Erklärung --- es justiert die Feinabstimmung empirisch
auf den Wert, der beobachtet wird. Supersymmetrie, Technicolor,
Randall-Sundrum-Modelle: alle sind Versuche, diesen Skalenabstand
mechanisch zu erklären. Keiner hat sich bisher experimentell bestätigt.

Die Frage, die dieses Dokument stellt --- wie koppeln die Skalen der
Materie strukturell miteinander? --- ist also nicht nur eine
FFGFT-spezifische Lücke. Sie ist eine \emph{allgemein ungelöste Frage
der theoretischen Physik}. FFGFT benennt einen Mechanismus, der diese
Frage angehen soll: die durchgehende $\xi$-Geometrie auf dem 4D-Torus,
die auf jeder Skala dieselbe Grundstruktur trägt, mit
skalenspezifischen Korrekturen über $K_{\mathrm{frak}}$. Ob dieser
Mechanismus die Antwort ist, ist offen. Dass die Frage offen ist, ist
es nicht.

\section{Die zwei Lese-Perspektiven}

Wenn man die Welt beschreiben will, gibt es zwei mögliche
Eintrittsperspektiven. Beide sind alt, beide haben in der Physik ihre
Tradition, und beide funktionieren auf ihre Weise.

\subsection*{Perspektive~1 --- die Bausteinsicht}

Die klassische reduktionistische Lesart: Welt besteht aus Quarks,
Quarks bauen Nukleonen, Nukleonen bauen Atomkerne, Kerne und Elektronen
bauen Atome, Atome bauen Moleküle, Moleküle bauen Kristalle oder Gase
oder Flüssigkeiten, diese bauen größere Strukturen --- Mineralien,
Gestein, Planetenschalen, Atmosphären, Hydrosphären --- und am Ende
Planeten, Sterne, Galaxien, Galaxienhaufen, das Universum als Ganzes.
Jede Stufe ist eine Ansammlung der vorhergehenden. Die Bindungskräfte
zwischen den Bausteinen werden lokal beschrieben: starke Wechselwirkung
auf der Kern-Skala, elektromagnetische auf der atomaren und
molekularen, van-der-Waals und Wasserstoffbrücken auf der
chemisch-strukturellen, Gravitation auf der astronomischen.

Die Standardphysik bewegt sich überwiegend in dieser Perspektive. Sie
beschreibt die Mechanismen \emph{innerhalb} jeder Stufe sehr gut. Sie
beschreibt auch die \emph{Übergänge} zwischen den Stufen ---
Festkörperphysik als Beschreibung, wie Atome zu Kristallen werden;
Hydrodynamik, wie Moleküle zu Flüssigkeiten werden; Astrophysik, wie
Materie zu Sternen wird. Was sie aber nicht beschreibt, ist, ob und
wie eine Geometrie auf einer Stufe \emph{durch} die anderen Stufen
\emph{hindurch} wirkt.

\subsection*{Perspektive~2 --- die Feldsicht}

Die holistische Lesart, die in der FFGFT zentral ist: Alle Stufen sind
\emph{stehende Wellen} desselben zugrunde liegenden Vakuumfeldes. Was
auf einer Stufe als ``Baustein'' erscheint, ist auf einer tieferen Ebene
eine stabilisierte Resonanzkonfiguration --- ein Elektron ist eine
Mode auf dem 4D-Torus, ein Atom ist eine zusammengesetzte Mode aus
Elektron- und Kern-Resonanzen, ein Kristall ist eine periodische Mode
über atomare Moden, ein Planet ist eine gravitativ stabilisierte
Resonanz aus Kristallaggregaten. Auf jeder Stufe gilt dieselbe
Geometrie, dieselbe Wicklungs-Logik, dieselbe Modenstruktur --- nur in
anderem Maßstab.

Die FFGFT-Aussage ist hier nicht, dass die Bausteinperspektive falsch
wäre. Sie ist nicht falsch. Aber sie ist eine \emph{lokale Sprache}, die
zwischen den Stufen Brüche hat, die in der Feldsicht keine Brüche sind.
Die Feldsicht hat eine andere Stärke: sie zeigt, warum die Stufen
\emph{gerade so} liegen, wie sie liegen --- nicht als zufällige Anhäufung
von Bausteinen, sondern als selbstkonsistente Resonanzkette.

\subsection*{Sind beide Perspektiven gleichwertig?}

Beide beschreiben dieselbe Welt. Sie sind nicht Konkurrenten, sondern
zwei Beschreibungssprachen für dasselbe. Welche ``richtiger'' ist, ist
eine missverstandene Frage --- so wenig wie ein Klavierstück richtiger
durch eine Notenpartitur oder durch ein akustisches Frequenzspektrum
beschrieben wird. Beide sind exakt; aber sie sehen verschiedenes
Verschiedenes.

\section{Die Skalenleiter}

Wenn man die Stufen ordnet --- vom kleinsten zum größten ---, ergibt
sich eine logarithmische Leiter, deren Sprossen ungefähr gleichmäßig
auseinander liegen. Jede Sprosse hat eine charakteristische Länge,
eine charakteristische Bindungsenergie, und einen logarithmischen
Abstand zur Planck-Skala.

\begin{center}
\small
\begin{tabular}{lcc}
\toprule
\textbf{Schicht} & \textbf{charakt.~Länge}~$\ell$ & $\log_{10}(\ell/\ell_{P})$ \\
\midrule
Planck-Skala       & $10^{-35}$~m & $0$ \\
Quark / Lepton     & $10^{-18}$~m & $\sim 17$ \\
Atomkern           & $10^{-15}$~m & $\sim 20$ \\
Atom               & $10^{-10}$~m & $\sim 25$ \\
Molekül            & $10^{-9}$~m  & $\sim 26$ \\
Kristall (Zelle)   & $10^{-9}$~m  & $\sim 26$ \\
Mikrostruktur (Korngrenze) & $10^{-6}$~m & $\sim 29$ \\
makroskop.~Materie & $10^{-3}$~m  & $\sim 32$ \\
Planetenkörper     & $10^{7}$~m   & $\sim 42$ \\
Stern              & $10^{9}$~m   & $\sim 44$ \\
Sonnensystem       & $10^{13}$~m  & $\sim 48$ \\
Galaxie            & $10^{21}$~m  & $\sim 56$ \\
Galaxienhaufen     & $10^{23}$~m  & $\sim 58$ \\
Universum (sichtbar) & $10^{26}$~m & $\sim 61$ \\
\bottomrule
\end{tabular}
\end{center}

Die Sprossen sind nicht gleichmäßig --- aber sie überspannen einen
logarithmischen Bereich von rund \emph{61 Größenordnungen}. Das ist die
volle Spannweite, über die die FFGFT-Geometrie wirken müsste, wenn die
Feldsicht durchgängig gelten soll.

\subsection*{Was die Standardphysik beschreibt}

Auf jeder Stufe gibt es eine etablierte Beschreibung:

\begin{itemize}\setlength\itemsep{2pt}
  \item \textbf{Quark/Lepton:} Standardmodell der Teilchenphysik (QCD,
        elektroschwach).
  \item \textbf{Atomkern:} Kernphysik (Schalenmodell, Mean-Field,
        ab-initio bei leichten Kernen).
  \item \textbf{Atom:} Quantenmechanik mit Coulomb-Wechselwirkung;
        Hartree-Fock, DFT.
  \item \textbf{Molekül:} Quantenchemie; ab-initio Methoden bis ca.
        100 Atome.
  \item \textbf{Kristall/Festkörper:} Festkörperphysik; Bändermodell,
        Bloch-Theorem, kollektive Anregungen (Phononen, Magnonen).
  \item \textbf{Flüssigkeit/Gas:} statistische Mechanik, Hydrodynamik,
        Thermodynamik.
  \item \textbf{Planetenkörper:} Geophysik, Mineralogie, klassische
        Mechanik.
  \item \textbf{Stern/Galaxie:} Astrophysik, Gravitationsphysik,
        Kosmologie.
\end{itemize}

Jede dieser Theorien ist auf ihrer Skala sehr genau. Aber jede arbeitet
mit einem \emph{eigenen Vokabular}, eigenen Konstanten, eigenen
Wechselwirkungen. Die Standardphysik gibt keine durchgehende Sprache,
die alle Skalen verbindet --- sie verschweißt sie durch Übergangs-Theorien.

\subsection*{Was FFGFT hinzufügen würde}

FFGFT würde sagen: auf \emph{jeder} dieser Stufen wirkt derselbe
Vakuum-Parameter $\xi_{0} = 4/30000$. Er wirkt überall mit derselben
lokalen Stärke (das ist die Aussage der Schicht~1, Dok.~241). Aber die
\emph{kumulative} Wirkung über den Skalenfluss zwischen den Stufen ist
verschieden, je nach logarithmischem Abstand. Dies ist die natürliche
Verallgemeinerung dessen, was in Dok.~133 für die elektroschwache
Skala explizit ausgearbeitet ist:
\[
  K_{\mathrm{frak}}(\Lambda_1 \to \Lambda_2) \;=\; 1 - \delta\!\cdot\!\xi\cdot\ln(\Lambda_1/\Lambda_2),
\]
mit $\delta$ als skalenspezifischem Koeffizienten, der von der
geometrischen Mittelung über die Stufen abhängt. Für den Sprung
Planck~$\to$~elektroschwach ist diese kumulierte Korrektur rund
$1{,}3\%$. Für Sprünge auf höheren Skalen sind die Korrekturen
\emph{kleiner}, nicht größer --- siehe nächster Abschnitt.

\section{Die zentrale Hypothese: kleine, aber durchgängige Kopplungen}

Hier kommt der zentrale Punkt dieses Dokuments, und er ist ehrlich als
\textbf{Hypothese, kein Beweis} ausgewiesen.

\subsection*{Aussage}

Die FFGFT-Geometrie wirkt auf allen Skalen, aber die direkten
geometrischen Kopplungen \emph{zwischen} weit voneinander entfernten
Skalen sind \emph{sehr klein}. Das ist eine notwendige Konsequenz der
Architektur:

\begin{enumerate}\setlength\itemsep{2pt}
  \item Der Parameter $\xi$ selbst ist klein ($\sim 10^{-4}$).
  \item Die Kopplung zwischen Skalen geht nur über den
        logarithmischen Abstand ein --- nicht über den linearen.
  \item Schon innerhalb des 17-Größenordnungs-Sprungs (Planck~$\to$
        elektroschwach) ist die kumulierte Korrektur nur $1{,}3\%$.
  \item Zwischen Skalen, die nur wenige Größenordnungen auseinander
        liegen (z.~B.\ Atom~$\to$~Kristall: $\sim 1$~Größenordnung),
        ist die direkte geometrische Kopplung um Faktoren $\sim 100$
        kleiner --- also $\sim 0{,}01\%$ oder darunter.
\end{enumerate}

Das bedeutet konkret: Eine Geometrie auf der atomaren Skala
``durchscheint'' in der makroskopischen Welt nur als Korrektur in der
Größenordnung von $10^{-4}$ bis $10^{-6}$ relativ zur dominanten
Wechselwirkung der jeweiligen Skala. Die Standardphysik ``sieht'' diese
Korrekturen meist nicht, weil sie unter ihrer Auflösungsgrenze liegen
und durch lokale Renormierungs- oder Mittelungsverfahren weggeräumt
werden.

\subsection*{Warum das die Bausteinperspektive rechtfertigt}

Genau weil die interskalaren Kopplungen klein sind, ist die
Bausteinperspektive der Standardphysik \emph{lokal sehr genau}. Wer
Atomphysik betreibt, kann ignorieren, dass dieselbe Geometrie auch in
einem Galaxienhaufen wirkt; die Beeinflussung liegt jenseits jeder
spektroskopischen Auflösung. Wer Kristallographie betreibt, kann
ignorieren, dass derselbe Vakuum-Parameter auch in der elektroschwachen
Theorie auftritt. Die Skalentrennung der Standardphysik ist nicht
falsch; sie ist eine sehr gute Näherung.

FFGFT widerspricht der Skalentrennung nicht. Sie macht nur eine
zusätzliche Aussage: die Trennung ist \emph{nicht absolut}. Auf
ausreichend hoher Präzision tauchen Spuren der durchgehenden Geometrie
auf --- als systematische Abweichungen, die sich nicht durch lokale
Mechanismen erklären lassen.

\subsection*{Wo solche Spuren beobachtbar sein könnten}

Folgende Bereiche sind die plausibelsten Kandidaten, an denen
interskalare FFGFT-Kopplungen experimentell sichtbar werden könnten,
wenn die Hypothese stimmt:

\begin{itemize}\setlength\itemsep{2pt}
  \item Hochpräzisions-Spektroskopie an Atomen und Molekülen
        (anomale Verschiebungen jenseits der QED-Korrekturen,
        Größenordnung $10^{-6}$ relativ).
  \item Anomale Korrekturen in Festkörper-Bandstrukturen, die sich
        nicht durch lokale Vielteilcheneffekte erklären lassen.
  \item Korrelationen zwischen makroskopischen Konstanten (etwa
        Materialdichten, kritischen Punkten, Phasenübergangs-Energien)
        und atomaren Größen, die über bekannte Skalierungsgesetze
        hinausgehen.
  \item Astronomische Verhältnisse (Planetenabstände, stellare
        Massen-Radius-Beziehungen), die sich systematisch leicht von
        rein klassisch-gravitativen Vorhersagen unterscheiden.
\end{itemize}

Alle vier Punkte sind als Forschungsrichtungen formuliert, nicht als
Vorhersagen. Niemand hat solche Kopplungen bisher systematisch in der
Größenordnung $10^{-4}$ bis $10^{-6}$ gesucht, weil keine etablierte
Theorie sie erwartet.

\section{Das mehrstöckige Konzertsaal-Bild}

Wer das Konzertsaal-Bild aus Dok.~159 kennt, kann hier eine
strukturelle Erweiterung sehen: \emph{der Konzertsaal hat mehrere
Stockwerke}.

Auf jedem Stockwerk gibt es einen eigenen akustischen Raum, mit eigenen
Resonanzen, eigener Nachhallzeit, eigener charakteristischer Frequenz.
Ein Ton, der im Erdgeschoss erzeugt wird, hat seine dominanten
Resonanzen im Erdgeschoss. Er dringt aber, durch Wände und Decken
hindurch, leise in das erste Stockwerk und in das zweite. Was im
Erdgeschoss ein voller Klang ist, ist im zweiten Stockwerk nur noch
ein schwaches Mitschwingen --- aber es ist nicht null. Und umgekehrt:
ein lauter Ton im obersten Stockwerk dringt leise hinunter.

So wird das mit den Skalen vorgestellt: Atomare Geometrie hat ihre
Resonanzen auf der atomaren Stufe, Festkörper-Geometrie auf der
Festkörper-Stufe, Planeten-Geometrie auf der Planeten-Stufe. Aber
jede Stufe schwingt leise in den anderen mit --- weil sie alle auf
demselben Vakuum stehen, in derselben fraktalen Geometrie. Die direkten
Kopplungen sind klein, weil zwischen weit auseinander liegenden Stufen
viele ``Wände und Decken'' dämpfen. Aber sie sind nicht null.

In dieser Sprache ist die Standardphysik diejenige Theorie, die jedes
Stockwerk \emph{einzeln} sehr gut beschreibt. FFGFT ist diejenige, die
fragt, wie das Gebäude als Ganzes klingt --- und die behauptet, dass
es leise Querkopplungen gibt, die ein einzelnes Stockwerk allein nicht
erklären kann.

\section{Die zwei Sprachen und ihre Synthese}

Damit lässt sich die anfangs formulierte Frage --- wie verbinden sich
Bausteinsicht und Feldsicht ohne Metapher? --- in folgender Form
beantworten:

\begin{itemize}\setlength\itemsep{2pt}
  \item Die \emph{Bausteinsicht} beschreibt jede Stufe \emph{lokal}
        und ihre direkten Übergänge zur nächst-höheren oder
        nächst-tieferen Stufe. Sie ist die Sprache der lokalen
        Mechanismen.
  \item Die \emph{Feldsicht} beschreibt das gemeinsame Vakuum, auf dem
        alle Stufen stehen, und die durchgängige Geometrie, die in
        allen Stufen dieselbe ist. Sie ist die Sprache der globalen
        Konsistenz.
  \item Beide sind exakt auf ihrer eigenen Auflösungsebene. Bausteinsicht
        ist exakter bei Beschreibung einer einzelnen Stufe. Feldsicht
        ist exakter bei der Frage, warum die Stufen so liegen, wie sie
        liegen, und wie sie selbstkonsistent miteinander verbunden sind.
  \item Die direkten interskalaren Kopplungen sind \emph{klein} --- so
        klein, dass die Bausteinsicht praktisch immer ausreicht. Aber
        sie sind nicht null --- und auf ausreichend hoher
        experimenteller Präzision sollten sie als systematische
        Korrekturen sichtbar werden.
\end{itemize}

\section{Was wir nicht sagen}

Wir sagen \emph{nicht}:

\begin{itemize}\setlength\itemsep{2pt}
  \item dass die Standardphysik falsch wäre;
  \item dass alle Mechanismen zwischen den Skalen aus FFGFT alleine
        ableitbar wären;
  \item dass das mehrstöckige Konzertsaal-Bild eine vollständige
        Theorie wäre;
  \item dass die Kopplungsstärken quantitativ aus der FFGFT bekannt
        sind --- sie sind \emph{nicht} aus den bisherigen Dokumenten
        ableitbar, sondern bisher nur als Größenordnungen plausibel.
\end{itemize}

Wir sagen, ehrlich als Hypothese:

\begin{itemize}\setlength\itemsep{2pt}
  \item dass die Geometrie des Vakuums auf allen Skalen wirkt;
  \item dass die Skalentrennung der Standardphysik eine sehr gute, aber
        \emph{nicht perfekte} Näherung ist;
  \item dass die interskalaren Kopplungen klein sind, aber prinzipiell
        in hochpräzisen Experimenten suchbar wären;
  \item dass dieser Aspekt der FFGFT ein offenes Forschungsterrain ist.
\end{itemize}

\section{Was zu tun wäre}

Um aus dieser Skizze eine prüfbare Theorie zu machen, wären drei
Schritte notwendig, die in zukünftigen Dokumenten anzugehen sind:

\begin{enumerate}\setlength\itemsep{2pt}
  \item \emph{Eine quantitative Form der Kopplungsformel} --- wie hängt
        die kumulierte Korrektur $K_{\mathrm{frak}}^{(n,m)}$ zwischen
        zwei Skalen $\Lambda_{n}$ und $\Lambda_{m}$ explizit vom
        logarithmischen Abstand ab? Dok.~133 enthält den Fall
        Planck~$\to$~elektroschwach; die Verallgemeinerung steht aus.
  \item \emph{Eine Auswahl konkreter Beobachtungsstellen}, an denen
        die interskalare Kopplung am besten testbar wäre, und die
        Vorhersage der erwarteten Effektgröße.
  \item \emph{Ein Vergleich mit der Standardphysik}, in dem ehrlich
        markiert wird, an welcher Stelle die FFGFT eine Vorhersage
        macht, die über die Standardmodell-Vorhersage hinausgeht ---
        und an welcher Stelle sie nur eine alternative Sprache für
        denselben Sachverhalt ist.
\end{enumerate}

Alle drei Schritte sind nicht-trivial. Sie sind das, was zwischen
einer Skizze wie diesem Dokument und einer Vollform der Theorie liegt.
Dieses Dokument macht keinen dieser Schritte; es markiert nur, dass
sie nötig sind.

\section*{Zusammenfassung}

\begin{itemize}\setlength\itemsep{2pt}
  \item FFGFT ist bisher am untersten Layer (Elementarteilchen) und
        am obersten (Kosmologie) ausgearbeitet; die Zwischenlayer
        --- Atome, Moleküle, Kristalle, Materie-Aggregate, Planeten,
        Sterne, Galaxien --- sind nur punktuell adressiert.
  \item Welt lässt sich in zwei komplementären Sprachen lesen:
        Bausteinsicht (lokal, Standardphysik) und Feldsicht (global,
        FFGFT). Beide sind exakt; sie sehen verschiedenes
        Verschiedenes.
  \item Die FFGFT-Geometrie wirkt nach Annahme auf allen Skalen mit
        derselben lokalen Stärke. Die kumulative Wirkung über
        Skalenabstände ist klein, weil die Abstände groß sind: Spuren
        in der Größenordnung $10^{-4}$ bis $10^{-6}$ relativ.
  \item Das rechtfertigt die Skalentrennung der Standardphysik als
        sehr gute Näherung --- nicht aber als prinzipielle Grenze.
  \item Status: \textbf{Hypothese, kein Beweis}. Quantitative
        Verallgemeinerung von $K_{\mathrm{frak}}$ auf beliebige
        Skalenpaare und konkrete Vorhersagen stehen aus.
\end{itemize}

\medskip
\noindent\textit{Stand: Mai 2026. Dieses Dokument ist
methodisch-pädagogisch und explizit als Plausibilitätsskizze
ausgewiesen. Es enthält keine eigenständigen quantitativen
Vorhersagen.}
